\chapter{Konstrukcja procesów i regularność trajektorii}
\label{ch:konstrukcja}

\section{Problem istnienia: od tabel rozkładów do całego procesu}
Załóżmy, że chcemy opisać losowy ruch bez podawania jawnego wzoru
$X_t(\omega)$. Możemy zamiast tego ustalić, jak mają wyglądać
rozkłady wszystkich wektorów
$(X_{t_1},\ldots,X_{t_m})$. Nie każda dowolna lista rozkładów
może jednak pochodzić od jednej rodziny zmiennych losowych.
Jeżeli np. zadamy $\Prob(X_1=0)=1$, a jednocześnie
$\Prob(X_1=1,X_2=0)=1$, otrzymamy sprzeczność.
Rozkłady muszą dawać zgodne odpowiedzi na pytania o te same
pomiarowe współrzędne.

W rozdziale~\ref{ch:procesy} wykazaliśmy, że rzeczywiste
rozkłady procesu spełniają relacje permutacji i marginalizacji.
Teraz udowodnimy twierdzenie odwrotne. Jego zasadnicza trudność
nie leży w sprawdzeniu skończenie wielu zmiennych, lecz w
zbudowaniu miary na przestrzeni funkcji indeksowanych
nieprzeliczalnym zbiorem chwil.

\begin{definicja}[Zgodny układ rozkładów]
Niech $T$ będzie niepustym zbiorem, a $E$ niepustą
przestrzenią polską, tj. ośrodkową przestrzenią topologiczną,
której topologię można zadać metryką zupełną.
Wyposażamy $E$ w borelowskie $\sigma$-ciało $\calB(E)$. Dla każdej skończonej
\emph{listy} $\mathbf t=(t_1,\ldots,t_m)$ indeksów z $T$
(niekoniecznie parami różnych) niech $\mu_{\mathbf t}$
będzie miarą probabilistyczną na $E^m$.
Rodzinę $(\mu_{\mathbf t})$ nazywamy \emph{zgodną}, jeżeli:
\begin{enumerate}
\item dla każdej permutacji $\pi$ indeksów $1,\ldots,m$,
$\mu_{t_{\pi(1)},\ldots,t_{\pi(m)}}=(\Pi_\pi)_\#\mu_{t_1,\ldots,t_m}$;
\item dla $k<m$ rzut $E^m\to E^k$ na pierwsze $k$ współrzędnych
przenosi $\mu_{t_1,\ldots,t_m}$ w $\mu_{t_1,\ldots,t_k}$;
\item jeżeli $t_i=t_j$, to $\mu_{t_1,\ldots,t_m}$ nadaje
prawdopodobieństwo $1$ zbiorowi $\{x\in E^m:x_i=x_j\}$.
\end{enumerate}
Symbol $(f)_\#\mu$ oznacza miarę obrazową
$((f)_\#\mu)(A)=\mu(f^{-1}(A))$.
\end{definicja}

Warunek dotyczący powtórzonych indeksów nie jest dodatkową
``własnością niezależności''. Wyraża po prostu to, że
$X_t$ oglądane dwa razy w tej samej chwili musi być tą samą
zmienną. Równoważnie można zadawać rozkłady tylko dla
skończonych \emph{zbiorów} czasów i wymagać zgodności
względem wszystkich rzutowań; wtedy problem powtórzeń znika.

\begin{lemat}[Algebra cylindrów]
Niech $E\ne\varnothing$ będzie przestrzenią polską, a $T$
niepustym zbiorem. Niech $\cal A$ będzie rodziną podzbiorów
$E^T$ postaci $\pi_J^{-1}(B)$, gdzie
$J\subseteq T$ jest skończony, $\pi_J:E^T\to E^J$
jest rzutowaniem, a $B\in\calB(E^J)$.
Wówczas $\cal A$ jest algebrą zbiorów i
$\sigma(\cal A)=\calB(E)^{\otimes T}$.
\end{lemat}
\begin{proof}
Zbiory $\varnothing,E^T$ należą do $\cal A$.
Dopełnienie $\pi_J^{-1}(B)$ jest równe
$\pi_J^{-1}(E^J\setminus B)$. Jeśli
$A=\pi_J^{-1}(B)$, $A'=\pi_K^{-1}(B')$, to
$A\cup A'$ jest cylindrem zależnym od skończonego zbioru
$J\cup K$: jego baza jest sumą przeciwobrazów $B,B'$ przez
odpowiednie rzuty $E^{J\cup K}\to E^J,E^K$.
Dowodzi to domknięcia na skończone sumy i dopełnienia.
Dla skończonego $J$ zachodzi
$\mathcal B(E^J)=\mathcal B(E)^{\otimes J}$.
Istotnie, rzuty są ciągłe, co daje jedną inkluzję,
a przeliczalna baza topologii $E$ daje przeliczalną
bazę prostokątów otwartych w $E^J$, co daje drugą inkluzję.
Każdy cylinder jest więc mierzalny względem produktowego
$\sigma$-ciała, a cylindry zależne od jednej współrzędnej
generują to $\sigma$-ciało. Stąd równość.
\end{proof}


\begin{przyklad}[Dlaczego nie wystarczy zgodność pojedynczych rozkładów?]
Niech $T=\{0,1\}$, $E=\{0,1\}$ i zażądajmy, aby
$X_0$ oraz $X_1$ miały symetryczny rozkład Bernoulliego.
Te wymagania dopuszczają zarówno niezależność
($\Prob(X_0=X_1)=1/2$), jak i równość $X_0=X_1$ prawie
na pewno. Trzeba więc wybrać także \emph{wspólne} prawo pary.
Jeśli rozszerzymy zbiór indeksów o kolejne chwile,
ten wybór musi być zgodny z prawami trójek, czwórek itd. Twierdzenie Kołmogorowa mówi, że przy
zgodności \emph{całej} rodziny problem jest rozwiązywalny;
nie mówi natomiast niczego o gładkości lub ciągłości
otrzymanych trajektorii.
\end{przyklad}

\begin{twierdzenie}[Twierdzenie Kołmogorowa o istnieniu procesu]
\label{thm:kol-ext}
Niech $T$ będzie dowolnym niepustym zbiorem, a $(E,d)$
niepustą przestrzenią polską z borelowskim $\sigma$-ciałem.
Niech dla każdej skończonej listy $\mathbf t=(t_1,\ldots,t_m)$
indeksów z $T$ dana będzie miara probabilistyczna
$\mu_{\mathbf t}$ na $E^m$, przy czym rodzina tych miar
jest zgodna w sensie powyższej definicji.
Wówczas istnieją przestrzeń probabilistyczna
$(\Omega,\calF,\Prob)$ i proces
$(X_t)_{t\in T}$ o stanach w $E$ takie, że
\[
\Law(X_{t_1},\ldots,X_{t_m})=\mu_{t_1,\ldots,t_m}
\quad\text{dla wszystkich }m\ge1\text{ i }t_1,\ldots,t_m\in T.
\]
Można wybrać $\Omega=E^T$,
$\calF=\calB(E)^{\otimes T}$ i $X_t(\omega)=\omega(t)$.
Prawo procesu na tej przestrzeni jest jednoznacznie wyznaczone
przez rodzinę $\mu_{\mathbf t}$.
\end{twierdzenie}

\begin{proof}
\emph{Krok 1: miara na cylindrach.}
Niech $\cal A$ będzie algebrą z poprzedniego lematu.
Dla cylindra $A=\pi_J^{-1}(B)$ definiujemy
$\mu_0(A)=\mu_J(B)$, gdzie dla skończonego $J$
przez $\mu_J$ rozumiemy rozkład we współrzędnych $J$
w dowolnej ustalonej kolejności, utożsamiony z miarą na $E^J$.
Dla $J=\varnothing$ przyjmujemy, że $E^J$ jest jednopunktowe
i $\mu_J$ jest jedyną miarą probabilistyczną na tej przestrzeni. Relacje permutacji i
marginalizacji dowodzą, że wynik nie zależy od kolejności
ani od dołączenia sztucznych współrzędnych. W szczególności
wartość $\mu_0(A)$ nie zależy od wybranej prezentacji
cylindra: jeżeli $A=\pi_J^{-1}(B)=\pi_K^{-1}(B')$,
obie prezentacje przenosimy na wspólny skończony
zbiór współrzędnych $L=J\cup K$.
Rzut $\pi_L:E^T\to E^L$ jest surjekcją, gdyż $E$ jest
niepuste, więc równość cylindrów oznacza równość ich
baz po przeniesieniu na $E^L$. Zgodność rozkładów daje
równość ich miar pod $\mu_L$.
Stąd $\mu_0(E^T)=1$ i $\mu_0$ jest skończenie addytywna:
dla rozłącznych cylindrów $A,B$ wyrażamy oba na
wspólnym skończonym zbiorze współrzędnych i stosujemy
addytywność miary skończeniewymiarowej.

\emph{Krok 2: dlaczego potrzebujemy ciągłości z góry.}
Aby rozszerzyć $\mu_0$ na $\sigma(\cal A)$ za pomocą
twierdzenia Carathéodory'ego, musimy pokazać przeliczalną
addytywność $\mu_0$ na $\cal A$. W przypadku skończonej
miary skończenie addytywnej wystarczy następujący warunek:
\[
A_n\in\cal A,\quad A_{n+1}\subseteq A_n,\quad
\bigcap_{n\ge1}A_n=\varnothing
\quad\Longrightarrow\quad \mu_0(A_n)\longrightarrow0.
\tag{$*$}
\]
Rzeczywiście, jeżeli $A=\bigsqcup_{k\ge1}B_k$,
przy $A,B_k\in\cal A$, to zbiory
$R_n=A\setminus\bigcup_{k=1}^n B_k$ należą do $\cal A$,
maleją do $\varnothing$, a skończona addytywność daje
$\mu_0(A)=\sum_{k=1}^n\mu_0(B_k)+\mu_0(R_n)$.
Przejście do granicy daje przeliczalną addytywność.

\emph{Krok 3: dowód warunku $(*)$ przy użyciu zwartości.}
Przypuśćmy, że $A_n\downarrow\varnothing$ i że
$\mu_0(A_n)\not\to0$. Ponieważ $\mu_0$ jest nieujemna i skończenie addytywna,
jest monotoniczna, więc ciąg $\mu_0(A_n)$ jest nierosnący. Istnieje $\varepsilon>0$ takie, że
$\mu_0(A_n)\ge\varepsilon$ dla wszystkich $n$.
Każdy $A_n$ zależy od pewnego skończonego zbioru
współrzędnych $J_n\subseteq T$. Ich suma
$S:=\bigcup_{n\ge1}J_n$ jest co najwyżej przeliczalna.
Jeśli $S=\varnothing$, wszystkie cylindry $A_n$ są
równe $\varnothing$ albo $E^T$, więc sprzeczność jest
natychmiastowa. W przeciwnym razie wypiszmy elementy
$S$ jako $s_1,s_2,\ldots$ (dla skończonego $S$
lista kończy się po skończenie wielu wyrazach). Ponieważ rozkłady
jednowymiarowe są miarami borelowskimi na przestrzeni
polskiej, są ciasne na mocy twierdzenia o ciasności
skończonych miar na przestrzeniach polskich
(por. skrypt \emph{Teoria miary i całki}, rozdział
\emph{Miary borelowskie, regularność i miary Radona}): dla każdego $j$ istnieje niepusty
zwarty $K_j\subseteq E$ taki, że
\[
\mu_{s_j}(E\setminus K_j)<\varepsilon2^{-j-3}.
\tag{1}
\]
Dla każdego $n$ zapiszmy $A_n=\pi_{J_n}^{-1}(B_n)$.
Skończony produkt przestrzeni polskich jest polski
(metryka będąca skończoną sumą metryk zupełnych jest
zupełna, a produkt skończenie wielu przeliczalnych
zbiorów gęstych jest przeliczalny i gęsty).
Wewnętrzna regularność względem zbiorów zwartych
skończonej miary borelowskiej $\mu_{J_n}$ na tej
przestrzeni daje niepusty zwarty
$C_n\subseteq B_n$ taki, że
\[
\mu_{J_n}(B_n\setminus C_n)<\varepsilon2^{-n-3}.
\tag{2}
\]
Zbiory $C_n$ są niepuste, ponieważ
$\mu_{J_n}(B_n)=\mu_0(A_n)\ge\varepsilon$.

Ustalmy $N$ i niech $I_N=\bigcup_{n=1}^N J_n$.
W skończonym produkcie $E^{I_N}$ rozważmy jednoczesne
warunki $\pi_{J_n}(x)\in C_n$ dla $1\le n\le N$
(tutaj $\pi_{J_n}$ oznacza rzut z $E^{I_N}$ na $E^{J_n}$)
oraz $x(s_j)\in K_j$ dla wszystkich $s_j\in I_N$.
Ich wspólna miara pod $\mu_{I_N}$ jest co najmniej
\[
\begin{aligned}
\mu_0(A_N)
&-\sum_{n=1}^N \mu_0\big(A_n\setminus\pi_{J_n}^{-1}(C_n)\big)
-\sum_{\{j:s_j\in I_N\}}\mu_{s_j}(E\setminus K_j)\\
&>\varepsilon-\frac\varepsilon8-\frac\varepsilon8
=\frac{3\varepsilon}{4}>0.
\end{aligned}
\tag{3}
\]
Użyliśmy tutaj $A_N\subseteq A_n$ dla $n\le N$ oraz
nierówności na miarę sumy, wyłącznie dla miary
$\mu_{I_N}$ (zgodność pozwala utożsamić jej wartości
na odpowiednich bazach cylindrów z wartościami $\mu_0$).
Nie korzystamy z przeliczalnej addytywności $\mu_0$,
której dopiero dowodzimy. Zatem istnieje wybór
współrzędnych na $I_N$, spełniający wszystkie te warunki.

Produkt $K:=\prod_{s_j\in S}K_j$ jest niepustą przestrzenią
zwartą (dla skończonego $S$ jest to produkt skończony).
Dla każdego $n$ zbiór
$F_n:=\{x\in K:\pi_{J_n}(x)\in C_n\}$
jest domknięty w $K$, ponieważ $C_n$ jest zwarty,
więc domknięty w $E^{J_n}$. Oszacowanie (3) pokazuje,
że dla każdego $N$ przecięcie $F_1\cap\cdots\cap F_N$
nie jest puste: znalezione wartości na skończonym $I_N$
można dopełnić dowolnymi punktami $K_j$ dla
pozostałych współrzędnych. Z własności zwartości
i skończonych przecięć dostajemy
$\bigcap_{n\ge1}F_n\ne\varnothing$.
Wybierzmy $x$ z tego przecięcia. Wówczas
$\pi_{J_n}(x)\in C_n\subseteq B_n$ dla każdego $n$.
Rozszerzając $x$ do elementu $E^T$ poprzez wybór
stałego punktu $e_0\in E$ poza $S$, otrzymujemy
$\omega\in\bigcap_{n\ge1}A_n$, sprzeczność.
Wobec tego warunek $(*)$ jest spełniony.

\emph{Krok 4: rozszerzenie i proces kanoniczny.}
Na mocy kroków 2--3 i twierdzenia Carathéodory'ego istnieje
miara probabilistyczna $\Prob$ na
$\calF=\sigma(\cal A)=\calB(E)^{\otimes T}$,
której ograniczenie do $\cal A$ jest równe $\mu_0$.
Odwzorowania współrzędne $X_t(\omega)=\omega(t)$ są
mierzalne z definicji produktowego $\sigma$-ciała,
a ich rozkłady dla list parami różnych indeksów są zadanymi
$\mu_{\mathbf t}$ z definicji $\mu_0$.
Dla dowolnej listy $\mathbf t=(t_1,\ldots,t_m)$ wybierzmy
listę $\mathbf s=(s_1,\ldots,s_k)$ wszystkich jej różnych
indeksów. Niech $r:E^m\to E^k$ wybiera po jednej
współrzędnej dla każdego $s_i$, a $j:E^k\to E^m$
odtwarza kolejność i powtórzenia z listy $\mathbf t$.
Z permutacji i marginalizacji wynika
$r_\#\mu_{\mathbf t}=\mu_{\mathbf s}$.
Na zbiorze, na którym wszystkie powtórzone współrzędne
są równe, zachodzi $j\circ r=\operatorname{id}$.
Zbiór ten jest mierzalny i ma miarę $1$ pod $\mu_{\mathbf t}$
na mocy trzeciego warunku zgodności. W konsekwencji
$\mu_{\mathbf t}=j_\#\mu_{\mathbf s}$.
Wektor kanonicznych współrzędnych dla listy $\mathbf t$
jest obrazem wektora dla $\mathbf s$ przez $j$, więc
również ma wymagany rozkład.
Jednoznaczność prawa wynika z jednoznaczności miar,
które zgadzają się na generującej algebrze cylindrów.
\end{proof}

\begin{uwaga}[Rola założenia o przestrzeni polskiej]
Założenie polskości $E$ pojawia się w kroku 3: gwarantuje,
że miary skończeniewymiarowe można od wewnątrz
przybliżać zbiorami zwartymi, a marginesy są ciasne.
Bez odpowiednich założeń o przestrzeni stanów
rozszerzenie zgodnej rodziny rozkładów może zawieść.
Twierdzenie nie zapewnia jeszcze ciągłości trajektorii.
Jeżeli zbiór indeksów jest wyposażony w $\sigma$-ciało,
nie gwarantuje też mierzalności wspólnej odwzorowania
$(t,\omega)\mapsto X_t(\omega)$ względem odpowiedniego
produktowego $\sigma$-ciała. Otrzymaliśmy tylko
mierzalność każdej współrzędnej $X_t$ z osobna.
\end{uwaga}

\section{Przykłady konstrukcji z zadanych rozkładów}
\begin{przyklad}[Nieskończona rodzina niezależnych zmiennych]
Niech $T$ będzie dowolnym zbiorem, a $\nu$ miarą
probabilistyczną na przestrzeni polskiej $E$.
Dla $t_1,\ldots,t_m$ parami różnych zadajmy
$\mu_{t_1,\ldots,t_m}=\nu^{\otimes m}$,
a dla powtórzonych indeksów utożsammy odpowiednie
współrzędne. Jest to rodzina zgodna, ponieważ
permutacja współrzędnych nie zmienia miary produktowej,
a rzut usuwa niepotrzebne czynniki.
Twierdzenie Kołmogorowa tworzy proces o niezależnych
$X_t$ i identycznych rozkładach $\nu$.

Jeżeli $T=[0,1]$ i $\nu$ jest rozkładem Bernoulliego z
parametrem $1/2$, to dla różnych $s,t$
$\Prob(|X_t-X_s|=1)=1/2$. Prawa strona nie dąży do
zera przy $s\to t$. Taki proces nie jest stochastycznie
ciągły i nie może mieć ciągłej modyfikacji: ciągłość
trajektorii prawie na pewno implikowałaby zbieżność
w prawdopodobieństwie, a modyfikacja zachowuje
rozkłady par $(X_s,X_t)$, a więc także rozkłady przyrostów. Samo
istnienie procesu jest więc znacznie słabsze niż
regularność jego trajektorii.
\end{przyklad}

\begin{przyklad}[Proces gaussowski z zadanej kowariancji]
Niech $T$ będzie zbiorem, $m:T\to\R$ dowolną funkcją,
a $K:T\times T\to\R$ funkcją symetryczną taką, że
dla wszystkich $k\ge1$, $t_1,\ldots,t_k\in T$
i $a_1,\ldots,a_k\in\R$ zachodzi
$\sum_{i,j=1}^k a_ia_jK(t_i,t_j)\ge0$.
Dla listy czasów $t_1,\ldots,t_k$ zadawajmy
$\mu_{t_1,\ldots,t_k}=\Normal(m_{\mathbf t},K_{\mathbf t})$,
gdzie $m_{\mathbf t}=(m(t_i))_{i=1}^k$,
$K_{\mathbf t}=[K(t_i,t_j)]_{i,j=1}^k$.
Nawet gdy $K_{\mathbf t}$ jest osobliwa, rozkład normalny
istnieje: można zapisać $m_{\mathbf t}+AZ$, gdzie
$Z\sim\Normal(0,I_k)$ i $AA^\top=K_{\mathbf t}$.
Macierz $A$ istnieje na mocy twierdzenia spektralnego
dla rzeczywistych macierzy symetrycznych: jeśli
$K_{\mathbf t}=Q\Lambda Q^\top$, gdzie $Q$ jest ortogonalna,
a $\Lambda$ diagonalna o nieujemnych wyrazach, można
wziąć $A=Q\Lambda^{1/2}$.
Rozkład ten jest jednoznacznie wyznaczony przez średnią
i kowariancję, gdyż jego funkcja charakterystyczna w punkcie
$u\in\R^k$ ma postać
$\exp(iu^\top m_{\mathbf t}-u^\top K_{\mathbf t}u/2)$
(por. rozdział o procesach gaussowskich).
Rzut i permutacja wektora gaussowskiego zachowują jego
gaussowskość i odpowiednio przekształcają średnią
oraz kowariancję. Z jednoznaczności rozkładu wynika
zatem zgodność z permutacją i marginalizacją. Jeżeli $t_i=t_j$, to różnica
odpowiednich współrzędnych ma średnią $m(t_i)-m(t_j)=0$
i wariancję $K(t_i,t_i)+K(t_j,t_j)-2K(t_i,t_j)=0$,
więc współrzędne te są równe prawie na pewno. Zatem rodzina jest zgodna
i istnieje proces gaussowski o średniej $m$ i kowariancji $K$.
\end{przyklad}


\begin{przyklad}[Konstrukcja łańcucha Markowa z macierzy przejścia]
Niech $E=\{0,1\}$, $T=\N_0$, $\mu=(\mu_0,\mu_1)$
będzie rozkładem początkowym, a
$P=(p_{ij})_{i,j\in E}$ macierzą stochastyczną:
$p_{ij}\ge0$ i $\sum_{j\in E}p_{ij}=1$ dla każdego $i$.
Zadajmy dla $n\ge0$ i $x_0,\ldots,x_n\in E$
\[
\mu_{0,1,\ldots,n}(\{(x_0,\ldots,x_n)\})
:=\mu_{x_0}\prod_{k=1}^n p_{x_{k-1}x_k}.
\tag{7}
\]
Sumując po $x_n$ i korzystając z
$\sum_{x_n\in E}p_{x_{n-1}x_n}=1$, widzimy,
że ostatnia współrzędna może zostać usunięta
bez zmiany prawa wcześniejszych współrzędnych.
Gdy pomijamy chwilę pośrednią $k$, pojawia się suma
$\sum_j p_{ij}p_{j\ell}=(P^2)_{i\ell}$,
a po pominięciu większej liczby chwil
odpowiednie potęgi $P$.
Dokładniej, dla $t_1,\ldots,t_m\in\N_0$ takich, że
$0\le t_1<\cdots<t_m$, i $x_1,\ldots,x_m\in E$ określamy
\[
\mu_{t_1,\ldots,t_m}(\{(x_1,\ldots,x_m)\})
=(\mu P^{t_1})_{x_1}
\prod_{k=2}^m (P^{t_k-t_{k-1}})_{x_{k-1}x_k}.
\]
Każda potęga $P^r$, $r\in\N_0$, jest macierzą stochastyczną:
iloczyn dwóch macierzy stochastycznych jest nieujemny
i zachowuje wektor kolumnowy jedynek. Sumowanie
powyższego wzoru po kolejnych współrzędnych od końca
daje zatem $1$, więc jest to miara probabilistyczna.
Równość $P^{a+b}=P^aP^b$ pokazuje zgodność
z dowolnymi rzutowaniami; powtórzenie chwili
oznacza z definicji powtórzenie tej samej współrzędnej.
Twierdzenie~\ref{thm:kol-ext} daje więc proces
$(X_n)_{n\ge0}$ o tych rozkładach.

Dla zdarzenia o dodatnim prawdopodobieństwie
$\{X_0=x_0,\ldots,X_n=x_n\}$ obliczamy z (7)
\[
\begin{aligned}
&\Prob(X_{n+1}=j\mid X_0=x_0,\ldots,X_n=x_n)\\
&\quad=\frac{\mu_{x_0}\,p_{x_0x_1}\cdots
p_{x_{n-1}x_n}\,p_{x_nj}}
{\mu_{x_0}\,p_{x_0x_1}\cdots p_{x_{n-1}x_n}}
=p_{x_nj}.
\end{aligned}
\]
Tak otrzymujemy własność Markowa, a nie tylko
jej formalną deklarację.

Dla konkretnego modelu weźmy
$\mu=(1,0)$ i
$P=\begin{pmatrix}3/4&1/4\\1/2&1/2\end{pmatrix}$.
Wtedy $\Prob(X_0=0,X_1=0,X_2=1)
=1\cdot(3/4)\cdot(1/4)=3/16$, podczas gdy
$\Prob(X_2=1)=(P^2)_{01}
=3/4\cdot1/4+1/4\cdot1/2=5/16$.
Różnica polega na tym, że w drugim pytaniu
sumujemy po wszystkich możliwych stanach w chwili $1$.
\end{przyklad}

\section{Kiedy proces może mieć trajektorie ciągłe?}
Rozkłady skończeniewymiarowe odpowiadają na pytania o
skończoną liczbę chwil, ale własność ciągłości wymaga
jednoczesnej kontroli \emph{nieskończenie wielu} przyrostów.
Narzędziem będzie oszacowanie momentów przyrostów
$\E\norm{X_t-X_s}^p$ o wykładniku potęgowym ściśle
większym niż $1$ w jednowymiarowym parametrze czasu.

\begin{definicja}[Ciągłość Höldera]
Niech $T>0$ oraz $\gamma>0$.
Funkcję $f:[0,T]\to\R^d$ nazywamy $\gamma$-h"olderowską,
jeśli istnieje stała $C_f<\infty$ taka, że
$\norm{f(t)-f(s)}\le C_f\abs{t-s}^{\gamma}$
dla wszystkich $s,t\in[0,T]$.
Gdy $\gamma>1$, funkcja o tej własności musi być stała:
dzieląc przedział $[s,t]$ na $n$ równych części,
to ten sam warunek daje
$\norm{f(t)-f(s)}\le C_f\abs{t-s}^{\gamma}n^{1-\gamma}\to0$.
\end{definicja}

\begin{lemat}[Kontrola przyrostów na siatkach dyadycznych]
\label{lem:dyadyczny}
Niech $T>0$, $d\ge1$ i $0<\gamma<\infty$.
Niech $D_n:=\{kT2^{-n}:k=0,\ldots,2^n\}$,
$D:=\bigcup_{n\ge0}D_n$. Jeżeli $f:D\to\R^d$
spełnia dla pewnego $N\ge0$
\[
\norm{f((k+1)T2^{-n})-f(kT2^{-n})}
\le 2^{-n\gamma}
\quad(n\ge N,\ 0\le k<2^n),
\]
to $f$ jest jednostajnie ciągła na $D$,
przedłuża się jednoznacznie do ciągłej funkcji
$\bar f:[0,T]\to\R^d$, a dla pewnej stałej
$C_{\gamma,T,N,f}<\infty$ zachodzi
$\norm{\bar f(t)-\bar f(s)}\le C_{\gamma,T,N,f}|t-s|^\gamma$.
\end{lemat}
\begin{proof}
Dla $u\in D$ i $n\ge N$ oznaczmy przez
$u_n=T2^{-n}\lfloor 2^n u/T\rfloor$ największy punkt
siatki $D_n$ nieprzekraczający $u$ (w szczególności
dla $u=T$ mamy $u_n=T$).
Dla dostatecznie dużego $m$ mamy $u_m=u$,
a kolejne przybliżenia $u_j,u_{j+1}$ są albo równe,
albo sąsiadują na siatce $D_{j+1}$. Z założenia
\[
\norm{f(u)-f(u_n)}
\le\sum_{j=n}^{m-1}2^{-(j+1)\gamma}
\le\frac{2^{-(n+1)\gamma}}{1-2^{-\gamma}}.
\tag{4}
\]
Jeśli $u,v\in D$ i $\abs{u-v}\le T2^{-n}$,
to ich lewostronne przybliżenia $u_n,v_n$
są identyczne lub sąsiadują na siatce $D_n$.
Wraz z (4) otrzymujemy
\[
\norm{f(u)-f(v)}\le
\left(1+\frac{2^{1-\gamma}}{1-2^{-\gamma}}\right)
2^{-n\gamma}.
\tag{5}
\]
Gdy $0<\abs{u-v}\le T2^{-N}$, wybierzmy $n\ge N$
tak, aby $T2^{-(n+1)}<\abs{u-v}\le T2^{-n}$.
Stąd $2^{-n\gamma}<(2\abs{u-v}/T)^\gamma$.
Z (5) dostajemy oszacowanie h"olderowskie na parach
wystarczająco bliskich. Dla par odległych o co najmniej
$T2^{-N}$ wykorzystujemy ograniczoność funkcji $f$ na
$D$, wynikającą z (4) oraz skończoności siatki $D_N$.
W ten sposób oszacowanie obowiązuje dla wszystkich par
w $D$ z pewną skończoną stałą. Ponieważ $\R^d$ jest
zupełna, $f$ ma jednoznaczne ciągłe przedłużenie
na domknięcie $D$, tj. na $[0,T]$; przejście do
odpowiednich granic przenosi oszacowanie na wszystkie $s,t$.
\end{proof}

\begin{twierdzenie}[Kryterium ciągłości Kołmogorowa]
\label{thm:kol-cont}
Niech $T>0$, $d\ge1$ i niech
$X=(X_t)_{t\in[0,T]}$ będzie procesem o wartościach w $\R^d$
na przestrzeni probabilistycznej $(\Omega,\calF,\Prob)$.
Załóżmy, że istnieją stałe $p>0$, $\beta>0$ i $C<\infty$
takie, że dla wszystkich $s,t\in[0,T]$ zachodzi
\[
\E\norm{X_t-X_s}^{p}\le C\abs{t-s}^{1+\beta}.
\tag{6}
\]
Wówczas istnieje modyfikacja $\widetilde X$ procesu $X$,
której trajektorie są ciągłe. Co więcej, można wybrać
$\widetilde X$ tak, że z prawdopodobieństwem $1$ jej
trajektorie są $\gamma$-h"olderowskie dla
\emph{każdego} $0<\gamma<\beta/p$.
\end{twierdzenie}
\begin{proof}
\emph{Krok 1: duże przyrosty na pojedynczej siatce.}
Ustalmy $0<\gamma_0<\beta/p$ i siatki
$D_n=\{kT2^{-n}:0\le k\le2^n\}$.
Niech $A_n$ będzie zdarzeniem, że istnieje sąsiednia
para punktów $D_n$, dla której przyrost przekracza
$2^{-n\gamma_0}$. Z nierówności Markowa, (6) i oszacowania
miary sumy mamy
\[
\begin{aligned}
\Prob(A_n)
&\le\sum_{k=0}^{2^n-1}
\frac{\E\norm{X_{(k+1)T2^{-n}}-X_{kT2^{-n}}}^{p}}
{2^{-np\gamma_0}}\\
&\le C T^{1+\beta}\,2^{-n(\beta-p\gamma_0)}.
\end{aligned}
\]
Ponieważ $\beta-p\gamma_0>0$, szereg
$\sum_{n\ge0}\Prob(A_n)$ jest zbieżny.

\emph{Krok 2: kontrola wszystkich dostatecznie drobnych siatek.}
Z pierwszego lematu Borela--Cantellego wynika, że z
prawdopodobieństwem $1$ zachodzi tylko skończenie wiele
zdarzeń $A_n$. Zatem na pewnym mierzalnym zbiorze
$\Omega_0$ o mierze $1$, dla każdego $\omega\in\Omega_0$
istnieje liczba całkowita $N(\omega)$ taka, że
wszystkie przyrosty sąsiednie na każdej siatce
$D_n$, $n\ge N(\omega)$, mają normę nie większą niż
$2^{-n\gamma_0}$. Dla takiego $\omega$ lemat
\ref{lem:dyadyczny} zastosowany do funkcji
$t\mapsto X_t(\omega)$ na $D:=\bigcup_nD_n$
daje jednoznaczne przedłużenie
$t\mapsto\widetilde X_t(\omega)$, które jest
$\gamma_0$-h"olderowskie.
Dla $\omega\notin\Omega_0$ ustawiamy
$\widetilde X_t(\omega)=0$ dla wszystkich $t$.
Przy ustalonym $t$ wybierzmy deterministyczny ciąg
$t_m\in D$ z $t_m\to t$. Na $\Omega_0$ mamy
$\widetilde X_t=\lim_m X_{t_m}$, a poza $\Omega_0$
zdefiniowaliśmy wartość stałą. Dokładniej, $\widetilde X_t$ jest granicą punktową na
całym $\Omega$ zmiennych mierzalnych
$\chi_{\Omega_0}X_{t_m}$: na $\Omega_0$ zachodzi zbieżność
do wartości przedłużenia, a poza nim wszystkie wyrazy
są zerowe. Nie wymaga to zupełności przestrzeni probabilistycznej.

\emph{Krok 3: przedłużenie jest modyfikacją.}
Dla ustalonego $t$ z (6) wynika
$\E\norm{X_{t_m}-X_t}^{p}\le C|t_m-t|^{1+\beta}\to0$.
Nierówność Markowa daje zbieżność $X_{t_m}\to X_t$
w prawdopodobieństwie, również gdy $0<p<1$.
Z drugiej strony $X_{t_m}\to\widetilde X_t$
prawie na pewno na $\Omega_0$, zatem także w
prawdopodobieństwie. Jednoznaczność granicy w
prawdopodobieństwie daje $\Prob(\widetilde X_t=X_t)=1$.
Ważne, że argument dotyczy osobno każdej ustalonej chwili;
nie zakładamy z góry nierozróżnialności.

\emph{Krok 4: wszystkie wykładniki mniejsze od $\beta/p$.}
Wybierzmy rosnący ciąg liczb
$0<\gamma_j<\beta/p$ taki, że
$\gamma_j\uparrow\beta/p$. Dla każdego $j$
powtórzmy kroki 1--2 z wykładnikiem $\gamma_j$.
Każdy z tych argumentów daje zbiór pełnej miary,
na którym ograniczenie procesu do \emph{tego samego}
zbioru $D$ ma ciągłe przedłużenie.
Przecinamy te przeliczalnie liczne mierzalne zbiory
pełnej miary również ze zbiorem $\Omega_0$ z kroku 2.
Na otrzymanym zbiorze jednoznaczność ciągłego
przedłużenia z $D$ zapewnia, że już skonstruowana
modyfikacja $\widetilde X$ ma trajektorie
$\gamma_j$-h"olderowskie dla wszystkich $j$.
Jeśli $0<\gamma<\beta/p$, wybieramy $j$ z
$\gamma<\gamma_j$ i wykorzystujemy
$|t-s|^{\gamma_j}\le T^{\gamma_j-\gamma}|t-s|^\gamma$.
To kończy dowód.
\end{proof}

\begin{uwaga}[Znaczenie wykładnika $1+\beta$]
Przy siatce dyadycznej liczba sąsiednich par jest rzędu
$2^n$, czyli rośnie odwrotnie proporcjonalnie do długości
kroku. Moment rzędu $p$ pojedynczego przyrostu jest
ograniczony z góry przez $CT^{1+\beta}2^{-n(1+\beta)}$. Po zsumowaniu po wszystkich parach
zyskujemy więc czynnik $2^{-n\beta}$.
Właśnie ta \emph{nadwyżka wykładnika ponad $1$}
pozwala zastosować Borela--Cantellego i otrzymać
jedno wspólne zdarzenie pełnej miary, kontrolujące
nieskończenie wiele chwil.
\end{uwaga}

\begin{przyklad}[Jak wybrać moment w kryterium Kołmogorowa?]
Załóżmy, że dla pewnego $p=4$ i $s,t\in[0,T]$
zachodzi $\E|X_t-X_s|^4\le C|t-s|^2$.
W zapisie twierdzenia mamy $1+\beta=2$, zatem $\beta=1$
i $\beta/p=1/4$. Otrzymujemy modyfikację o trajektoriach
h"olderowskich każdego rzędu $\gamma<1/4$.
Nie twierdzimy, że wykładnik $1/4$ jest optymalny:
może się okazać, że dysponujemy mocniejszymi
oszacowaniami dla wyższych momentów, prowadzącymi do
większych wykładników.
\end{przyklad}

\section{Konstrukcja procesu Wienera: pierwszy kontakt}
Rozdziały poświęcone ruchowi Browna i mierze Wienera
rozwiną tę konstrukcję, ale teraz możemy ją już wykonać
matematycznie poprawnie, korzystając z dwóch udowodnionych
twierdzeń. Zauważmy, że nie potrzebujemy żadnej całki
stochastycznej.

\begin{twierdzenie}[Istnienie standardowego procesu Wienera]
\label{thm:wiener-exists}
Istnieją przestrzeń probabilistyczna i proces rzeczywisty
$W=(W_t)_{t\ge0}$ takie, że:
\begin{enumerate}
\item $W_0=0$ prawie na pewno;
\item $W$ ma niezależne przyrosty;
\item dla $0\le s<t$ zachodzi
$W_t-W_s\sim\Normal(0,t-s)$;
\item trajektorie $t\mapsto W_t$ są ciągłe prawie na pewno.
\end{enumerate}
Można ponadto wybrać wspólny zbiór pełnej miary,
na którym trajektorie są $\gamma$-h"olderowskie
na każdym przedziale $[0,T]$ dla każdego $0<\gamma<1/2$.
\end{twierdzenie}
\begin{proof}
\emph{Krok 1: zgodne rozkłady gaussowskie.}
Dla $s,t\ge0$ niech $K(s,t)=\min(s,t)$.
Jest to funkcja dodatnio półokreślona, ponieważ
\[
\sum_{i,j=1}^m a_ia_j\min(t_i,t_j)
=\int_0^\infty
\left(\sum_{i=1}^m a_i\chi_{[0,t_i]}(u)\right)^2
\,\mathrm d\lambda_1(u)\ge0.
\]
Z przykładu o procesach gaussowskich i twierdzenia
\ref{thm:kol-ext} istnieje proces gaussowski $X$
o zerowej średniej i kowariancji $K$.
Ponieważ $K(0,0)=0$, $X_0=0$ prawie na pewno.

\emph{Krok 2: niezależne przyrosty i ich rozkłady.}
Dla $0\le s<t$ wariancja przyrostu wynosi
\[
\Var(X_t-X_s)=K(t,t)+K(s,s)-2K(s,t)=t-s.
\]
Jest to zmienna gaussowska o średniej $0$.
Dla dwóch rozłącznych przedziałów $(a,b]$ i $(c,d]$,
gdzie $a<b\le c<d$, obliczamy
\[
\Cov(X_b-X_a,X_d-X_c)
=K(b,d)-K(b,c)-K(a,d)+K(a,c)
=b-b-a+a=0.
\]
Wektor wszystkich przyrostów jest gaussowski, ponieważ
jest liniowym obrazem skończonego wektora gaussowskiego.
Jego macierz kowariancji jest diagonalna dla
parami rozłącznych przedziałów, więc funkcja charakterystyczna
wektora jest iloczynem funkcji charakterystycznych jego
współrzędnych. Z jednoznaczności funkcji charakterystycznych
wynika, że te przyrosty są \emph{łącznie niezależne}; sama parami zerowa kowariancja
nie wystarczyłaby dla dowolnych rozkładów, lecz wystarcza
dla wektora gaussowskiego.

\emph{Krok 3: ciągła modyfikacja.}
Jeśli $Z\sim\Normal(0,1)$, to dla każdego $p>0$
stała $c_p:=\E|Z|^p$ jest skończona, gdyż
$|z|^p\exp(-z^2/2)$ jest całkowalne na $\R$
(w nieskończoności wykładniczy zanik dominuje wzrost
potęgowy, a w otoczeniu zera funkcja jest ograniczona).
Z rozkładu przyrostów wynika
\[
\E|X_t-X_s|^p=c_p|t-s|^{p/2}.
\]
Dla $p>2$ wykładnik $p/2$ jest większy niż $1$.
Kryterium \ref{thm:kol-cont} na $[0,T]$ daje
ciągłą modyfikację o trajektoriach h"olderowskich
każdego rzędu $\gamma<1/2-1/p$.
Aby otrzymać jeden proces na $[0,\infty)$, stosujemy
konstrukcję na przedziałach $[0,n]$, $n\in\N$.
Na wspólnym zbiorze pełnej miary przedłużenia
zgadzają się na wspólnych punktach wymiernych,
bo wszystkie są modyfikacjami tego samego $X$.
Z ciągłości zgadzają się więc na całych częściach wspólnych
przedziałów. Oznaczmy ten mierzalny zbiór pełnej miary przez $G$,
a skonstruowane modyfikacje na $[0,n]$ przez $X^{(n)}$.
Dla $\omega\in G$ oraz $t\ge0$ definiujemy
$W_t(\omega)=X_t^{(n)}(\omega)$ dla dowolnego całkowitego
$n\ge\max(1,t)$, a dla $\omega\notin G$ kładziemy
$W_t(\omega)=0$. Zgodność na częściach wspólnych
zapewnia niezależność wyboru $n$ i ciągłość wszystkich
trajektorii. Dla ustalonego $t$ można wybrać deterministyczne
$n$, więc $W_t$ jest mierzalna; ponieważ $G$ ma miarę $1$
i $X_t^{(n)}=X_t$ prawie na pewno, $W$ jest modyfikacją $X$.

\emph{Krok 4: zachowanie rozkładów i wykładniki.}
Modyfikacja nie zmienia żadnego rozkładu
skończeniewymiarowego: dla skończonej liczby chwil
przecięcie zbiorów, na których $W_t=X_t$, ma miarę $1$.
Stąd $W$ zachowuje wszystkie własności z kroków 1--2.
Dla $n\ge1$ i $k\ge1$ wybierzmy $p_k=k+2$ i zastosujmy
kryterium ciągłości do ograniczenia $X$ na $[0,n]$.
Otrzymana modyfikacja jest nierozróżnialna z ograniczeniem
$W$ do $[0,n]$, na mocy twierdzenia o ciągłych
modyfikacjach z pierwszego rozdziału. Dla każdej pary
$(n,k)$ istnieje więc mierzalny zbiór pełnej miary,
na którym trajektoria $W$ jest $\gamma$-h"olderowska
na $[0,n]$ dla wszystkich $0<\gamma<1/2-1/p_k$.
Przecięcie tych przeliczalnie licznych zbiorów ma miarę $1$.
Dla dowolnych $T>0$ i $0<\gamma<1/2$ wybieramy
$n\ge T$ oraz $k$ tak duże, że $\gamma<1/2-1/p_k$,
i ograniczamy odpowiednie oszacowanie do $[0,T]$.
\end{proof}

\begin{uwaga}[Co dokładnie skonstruowaliśmy?]
Proces Wienera jest tutaj \emph{wnioskiem} z twierdzenia
o istnieniu rozkładów i kryterium ciągłości.
Jego pełną teorię -- zasadę odbicia, czasy trafienia,
własność Markowa, nieróżniczkowalność, wariację
kwadratową, most Browna i miarę Wienera --
przedstawimy w kolejnych rozdziałach tomu II.
Żadna z tych dodatkowych własności nie została
jeszcze użyta w konstrukcji.
\end{uwaga}

\begin{center}
\begin{tikzpicture}[>=Latex,scale=1]
\draw[->] (0,0)--(10.4,0) node[right] {$t$};
\draw[->] (0,-1.7)--(0,1.9) node[above] {$W_t$};
\draw[gray!40,step=0.5] (0,-1.5) grid (10,1.5);
\draw[very thick,blue!65!black]
(0,0)--(0.45,0.35)--(0.85,0.06)--(1.1,0.43)--(1.45,0.11)--
(1.8,0.57)--(2.15,0.30)--(2.35,0.74)--(2.65,0.28)--(3,0.46)--
(3.3,-0.05)--(3.62,0.37)--(3.85,0.04)--(4.1,-0.50)--(4.32,-0.22)--
(4.58,-0.84)--(4.82,-0.50)--(5.10,-0.95)--(5.34,-0.63)--(5.62,-1.12)--
(5.88,-0.61)--(6.12,-0.74)--(6.36,-0.25)--(6.60,-0.55)--(6.92,0.05)--
(7.20,-0.30)--(7.46,0.19)--(7.74,-0.12)--(8.02,0.30)--(8.28,0.02)--
(8.54,0.42)--(8.81,0.05)--(9.1,0.53)--(9.42,0.31)--(9.75,0.81);
\end{tikzpicture}
\par\smallskip\small Poglądowa łamana, a nie numerycznie wygenerowana realizacja procesu Wienera.
\end{center}

\section{Zadanie rozwiązane: dlaczego ciągłość wymaga osobnego dowodu?}
\begin{zadanie}[Proces o jednym losowym skoku]
Niech $U\sim\operatorname{U}([0,1])$ i
$X_t=\chi_{\{U\le t\}}$ dla $t\in[0,1]$.
\emph{(a)} Obliczyć $\Prob(X_s=0,X_t=1)$ dla $0\le s<t\le1$.
\emph{(b)} Wykazać stochastyczną ciągłość.
\emph{(c)} Sprawdzić, czy kryterium Kołmogorowa może
zostać zastosowane do procesu $X$.
\emph{(d)} Ustalić, czy istnieje ciągła modyfikacja.

\emph{Rozwiązanie.}
(a) Zdarzenie $\{X_s=0,X_t=1\}$ jest dokładnie
$\{s<U\le t\}$, a zatem ma prawdopodobieństwo $t-s$.
Podobnie $\Prob(X_s=1,X_t=0)=0$, co odzwierciedla
niemalejące trajektorie.

(b) Dla $s<t$ mamy $|X_t-X_s|\in\{0,1\}$ i
$\Prob(|X_t-X_s|=1)=t-s$. W konsekwencji dla każdego
$0<\varepsilon<1$ zachodzi
$\Prob(|X_t-X_s|>\varepsilon)=|t-s|\to0$.
Dla $\varepsilon\ge1$ prawdopodobieństwo wynosi $0$.
To dowodzi ciągłości stochastycznej.

(c) Dla każdego $p>0$ zachodzi
$\E|X_t-X_s|^p=|t-s|$, ponieważ zmienna
$|X_t-X_s|$ przyjmuje tylko wartości $0,1$.
Nie istnieją więc $C<\infty$ i $\beta>0$ takie, że
$|t-s|\le C|t-s|^{1+\beta}$ dla wszystkich $s,t$:
dzieląc przez $|t-s|>0$ i posyłając $t\to s$,
dostalibyśmy $1\le0$. Założenie kryterium
Kołmogorowa nie jest spełnione.

(d) Dla każdej realizacji z $U\in(0,1)$ funkcja
$t\mapsto X_t$ ma skok. Sam ten fakt nie wystarcza,
by odrzucić możliwość \emph{innej} modyfikacji.
Jednak każda modyfikacja musi zgadzać się z $X$
we wszystkich chwilach wymiernych poza jednym zbiorem
zerowym. Po obu stronach $U$ wartości w chwilach
wymiernych są odpowiednio równe $0$ i $1$; nie
mogą więc być wartościami funkcji ciągłej.
Ciągła modyfikacja nie istnieje.
\end{zadanie}

